\documentclass[10pt]{article}

\usepackage[margin=1in]{geometry}
\usepackage{amsmath,amssymb}
\usepackage{booktabs}
\usepackage{graphicx}
\usepackage{hyperref}
\usepackage{natbib}

\title{Belief-Guided Arbitration for Interactive Troubleshooting}
\author{Hanshal Goyal}
\date{}

\begin{document}
\maketitle

\begin{abstract}
Large language model agents often need to gather information before acting, but
directly transferring Bayesian experimental design (BED) to open-ended
interaction creates a proposal problem: a statistically informative action can
be a poor task action. We study this failure on the externally released Paprika
customer-service benchmark. A faithful BED-LLM transfer and full two-step
lookahead fail to improve over a thinking agent in a terminal-audited
development study. We therefore separate proposal from selection. A native
thinking policy generates three credible next actions; a belief model estimates
their one-step expected information gain (EIG); and a calibrated arbiter
overrides the native default only when an alternative clears a one-standard-error
margin. A ten-task development pilot motivates a preregistered 50-task held-out
evaluation against thinking and non-thinking agents and a prompt-matched
candidate-0 control. This draft fixes the method, benchmark adaptation,
measurement audit, and analysis protocol before held-out outcomes are opened.
The study tests whether explicit beliefs are more useful for selecting among
good native proposals than for generating actions from scratch.
\end{abstract}

\section{Introduction}

Interactive agents rarely receive all information needed for a correct action.
A support agent must identify a fault before recommending a remedy; a clinical
agent must elicit missing symptoms before diagnosing. Language models can ask
fluent questions, yet fluency does not guarantee useful information gathering.
MediQ, for example, finds that directly prompting models to ask questions can
reduce diagnostic performance \citep{li2024mediq}. The design problem is not
only what to say, but which uncertainty an action should resolve and whether
that information is worth a scarce interaction turn.

Bayesian experimental design gives a principled answer. Given hypotheses
$h\in\mathcal H$, posterior $p_t(h)$, and possible outcomes $y$, one-step EIG
selects the action expected to reduce posterior uncertainty. BED-LLM extends
this idea to semantic hypothesis spaces by using language models to propose
hypotheses and score likelihoods \citep{choudhury2025bedllm}. The appeal is
strong: a training-free scaffold can impose explicit information-seeking
structure on a general-purpose model.

Our first result in Paprika is a design diagnosis rather than a positive
transfer claim. A faithful EIG policy generated schema-valid diagnostic and
corrective actions but did not beat a thinking native policy in a terminal-valid
ten-task development study. Full two-step lookahead also failed. Inspecting the
interface exposed an uncontrolled distinction: the native policy was prompted
to solve the customer's issue, whereas the generic EIG generator was prompted
to produce valid candidates. EIG then ranked a weaker proposal set. More
lookahead could not repair missing task intent.

We respond by assigning the language model and BED scaffold different jobs. The
native policy proposes its natural next action and two credible alternatives.
The belief scaffold evaluates those same actions. Candidate 0 is the native
default; the arbiter selects an alternative only when its EIG advantage exceeds
a combined one-standard-error threshold. This rule treats score differences
inside estimated noise as a reason to retain the native action, not as evidence
for an arbitrary argmax.

The resulting study asks a narrow causal question: can explicit beliefs improve
selection among already-good native proposals? It includes a mandatory
candidate-0 control that receives the identical three-proposal prompt but always
executes the default. This distinguishes value from proposal elicitation from
value from EIG selection. Our contributions are:

\begin{enumerate}
  \item an external Paprika adaptation with explicit semantic beliefs,
  categorical answer spaces, likelihood-based EIG, and strict endpoint audits;
  \item a calibrated native-primary arbitration rule that combines a thinking
  policy's proposal quality with BED selection;
  \item a prompt-matched candidate-0 control and shared-candidate protocol that
  isolate the causal effect of EIG-based overrides; and
  \item a preregistered paired evaluation with whole-task recovery, outcome
  blindness, bootstrap uncertainty, and transcript-level endpoint review.
\end{enumerate}

\section{Related Work}

\paragraph{Bayesian experimental design.}
BED chooses actions by expected utility under a probabilistic model. Deep
adaptive design amortizes sequential policies and evaluates total information
gain with contrastive bounds \citep{foster2021deep}. BED-LLM uses language
models to extend hypothesis generation and likelihood elicitation to semantic
domains \citep{choudhury2025bedllm}. Our method keeps its one-step categorical
EIG foundation but changes the proposal-selection interface after faithful
end-to-end transfer failed in development.

\paragraph{Curiosity and interactive agents.}
Paprika trains a generally curious agent with reinforcement learning across
interactive environments \citep{tajwar2025curious}. We instead ask whether an
inference-time Bayesian scaffold can improve a fixed model without policy
training. MediQ turns medical question answering into an adaptive interaction
and shows that naive question asking is not reliably beneficial
\citep{li2024mediq}. These benchmarks motivate evaluating information gathering
through task endpoints rather than through question quality alone.

\paragraph{Planning and selective computation.}
Multi-step BED can value prerequisite questions whose benefit appears only in
future decisions, but branching over semantic outcomes is expensive and noisy.
Our development study tested full two-step lookahead and found no advantage, so
the held-out method deliberately returns to one-step EIG. The selectivity here
is statistical rather than temporal: the method selectively trusts an EIG
override only when the estimated score gap clears a fixed uncertainty margin.

\section{Paprika Troubleshooting Environment}

\subsection{External Task Definition}

We use the hash-verified official Paprika customer-service release. Each task
contains a public customer scenario and a private cause/remedy used by the
released customer simulator and success protocol. The private solution is never
available to proposal generation, belief construction, likelihood scoring, or
selection. It enters only the simulator and endpoint checks. This preserves the
benchmark's task structure instead of inventing a synthetic hidden state.

An interaction consists of up to five atomic diagnostic questions or corrective
actions. The customer replies in free text. A task resolves when the complete
conversation establishes or performs the specific private remedy and the
customer confirms success. Unresolved tasks are censored at six turns for the
primary turns-to-resolution endpoint.

\subsection{Beliefs, Outcomes, and Likelihoods}

For each public scenario, the questioner proposes a finite support of plausible
cause-and-remedy hypotheses. After every observation it proposes refinements,
filters only hypotheses contradicted by the public history, and recomputes the
posterior over the retained support. The implementation caps support size but
does not give the model access to the benchmark solution.

Free-text customer replies require an explicit answer-space interface. For each
candidate action $a$, the model proposes three to five mutually exclusive,
customer-observable outcomes. Every corrective action includes resolution,
attempted-but-unchanged, and unable-to-perform outcomes. For hypothesis $h$, an
LLM likelihood call estimates $p(y\mid h,a)$. A separate non-thinking mapper
maps the realized reply to one proposed outcome or marks the map unclean.
Answer-space coverage is the fraction of realized replies mapped cleanly.

The one-step action score is
\[
  I_t(a)=\sum_h p_t(h)\sum_y p(y\mid h,a)
  \log\frac{p(y\mid h,a)}{\sum_{h'}p_t(h')p(y\mid h',a)}.
\]
All candidate actions at a decision share the same posterior support.

\subsection{Endpoint Validation}

Open-ended simulators can invalidate an otherwise careful policy comparison. A
simulator may claim that a plausible but incorrect remedy succeeded, or claim
that the private remedy failed. We therefore use a two-layer check. First, each
reply is audited against the private solution and regenerated within a bounded
repair budget if inconsistent. Second, every terminal ``Goal reached'' claim is
checked against the latest action and private remedy. Surviving contradictions
fail the complete task invocation closed.

The adapter also runs Paprika's complete-conversation success judge after every
turn. Before policy evaluation, a five-task smoke required at least 85\% clean
answer-space coverage, zero surviving simulator contradictions, zero structured
failures, and manual transcript review. An earlier endpoint implementation was
quarantined when manual audit found false successes; none of those outcomes are
used as policy evidence. Endpoint auditing is therefore part of the method, not
post-hoc cleanup.

\section{Belief-Guided Native Arbitration}

At each turn, the thinking questioner sees the public scenario and conversation
and returns exactly three ordered, atomic actions. Candidate 0 is the action it
would naturally take next. Candidates 1 and 2 are distinct credible
alternatives. The same prompt also supplies a discrete observable outcome set
for each action. The arbiter scores all three actions under the current belief.

To quantify score uncertainty, let
\[
  c_h(a)=\sum_y p(y\mid h,a)
  \log\frac{p(y\mid h,a)}{p(y\mid a)}
\]
be the expected information contribution under hypothesis $h$. We compute a
weighted variance of $c_h(a)$ under $p_t(h)$ and divide by the posterior
effective sample size $n_{\mathrm{eff}}=1/\sum_h p_t(h)^2$. Its square root is
the action's standard error $s(a)$.

Let $a_0$ be the native default and $a_*=\arg\max_a I_t(a)$. The deployed rule is
\[
  a_t=\begin{cases}
  a_*, & I_t(a_*)-I_t(a_0)>
  \sqrt{s(a_*)^2+s(a_0)^2},\\
  a_0, & \text{otherwise.}
  \end{cases}
\]
The strict inequality and one-SE threshold were frozen after the development
pilot. The method logs all candidates, scores, standard errors, margins, and
override decisions.

\paragraph{Candidate-0 control.}
The causal control receives the identical ordered three-proposal prompt and
always executes $a_0$. It constructs no beliefs and performs no EIG scoring.
Whenever its public history matches arbitration, a shared prompt cache makes
the ordered proposal set byte-for-byte identical after parsing. The held-out
analyzer rejects any mismatch. This control separates the effect of eliciting
three goal-directed actions from the effect of selecting among them with EIG.

\paragraph{Best-$N$ elicitation probe.}
A preregistered development-only probe tests a second explanation for faithful
EIG's failure. It changes only the generic candidate instruction to request the
five best next actions for resolving the issue quickly, then applies ordinary
EIG argmax. Tasks and outcomes are disjoint from the held-out evaluation. If
this variant matches arbitration, the simpler interpretation is goal-anchored
generation plus EIG selection; if it remains weaker, calibrated native-primary
selection is load-bearing.

\section{Evaluation Protocol}

The held-out Paprika study contains the next 50 unseen official evaluation
tasks, offsets 10--59 in released order, with seed 1304 and a five-turn budget.
The thinking comparison jointly evaluates arbitration, candidate 0, and a
native thinking agent. A separate non-thinking native agent provides context.
All arms use the same 26B A4B model route; thinking is enabled only for policy
generation in the designated arms. Simulator, mapping, likelihood, filtering,
and endpoint calls are non-thinking.

The primary endpoint is paired censored turns to resolution for arbitration
minus thinking naive. The co-primary causal endpoint compares arbitration with
candidate 0. We report mean paired differences with 5,000-resample bootstrap
95\% intervals, win/loss/tie counts, and Wilcoxon signed-rank tests. Resolution
at five turns uses paired discordance counts and an exact two-sided binary test.
The method claim requires both co-primary bootstrap intervals to exclude zero
in arbitration's favor.

The mechanism analysis reports override rate, immediate resolution after an
override, and score-margin distributions for arbitration wins, losses, and ties
against candidate 0. These analyses are descriptive; the threshold cannot be
retuned. Cost reporting includes requests, tokens, model cost, and cost per
resolution. Candidate-0 cost is labeled cache-amortized because shared proposals
can be reused while histories match.

Outcome blindness is enforced during execution. Monitoring can inspect process
health, structured failures, candidate pairing, and cost, but not resolution or
turn outcomes. A failed thinking unit discards all three methods for that task
and replays the complete unchanged triplet; partial methods are never combined.
The final analyzer runs once after all 50 canonical tasks and the non-thinking
blocks are fixed.

After automated analysis, manual review covers every task where arms disagree
on success plus ten seed-selected remaining tasks. Reviewers compare each
terminal claim and near-remedy against the private benchmark remedy. Any
surviving false success or false failure quarantines the complete headline
result rather than dropping a convenient task.

\section{Results Status}

The terminal-valid ten-task development pilot selected the arbitration rule for
held-out testing: arbitration resolved 60\% of tasks versus 40\% for thinking
naive, with six paired turn wins, no losses, four ties, and a mean censored-turn
difference of $-1.2$ (development bootstrap interval $[-2.2,-0.4]$). It
overrode the native default on 12 of 33 decisions. These values are explicitly
development evidence and do not establish the paper's method claim.

The 50-task held-out outcomes remain sealed while this methods draft is written.
No result-dependent wording, subgroup, threshold, endpoint, or task is selected
from partial runs. The final paper will populate this section only from the
frozen analyzer and mandatory manual endpoint audit.

\section{Discussion}

The central hypothesis is an interface claim. In open semantic tasks, native
LLM competence may be most valuable for constructing a small set of actions
that are actually useful, while explicit Bayesian structure may be most
valuable for discriminating among them. This differs from asking EIG to both
invent and select actions. The candidate-0 control makes that distinction
testable rather than rhetorical.

The calibrated margin also changes what an EIG score means operationally. An
argmax always returns a winner, even when candidate differences are negligible
relative to support variability. Native-primary arbitration instead treats
small score gaps as insufficient evidence to override a capable policy. This
does not prove the standard-error estimate is calibrated in a frequentist
sense; the held-out mechanism analysis only asks whether the frozen rule yields
useful overrides.

The broader target is transfer to MediQ, where hypothesis tracking corresponds
to differential diagnosis and questions elicit missing clinical facts. That
experiment is conditional on Paprika's held-out and manual gates. It is not part
of the present evidence and will not be launched to rescue a null Paprika
result.

\section{Limitations}

The study uses one model family and one externally released troubleshooting
benchmark. The 50-task sample is workshop scale and has approximately 78\%
planned power for half the development effect, slightly below the conventional
80\% target. LLM likelihoods and finite generated hypothesis supports are not
calibrated probabilistic models; the one-SE margin measures variation across
that support, not all epistemic uncertainty.

The customer simulator, mapper, likelihood model, and endpoint judges share a
model family, which can induce correlated errors despite role separation and
private-solution audits. The strict terminal protocol reduces false endpoint
claims but may fail closed on otherwise usable interactions, requiring exact
whole-task replay. Manual review is therefore essential. OpenRouter execution
also introduces provider nondeterminism despite fixed prompts and seeds.

Candidate 0 is a prompt-matched native default, not an independently sampled
copy of the plain naive policy. This is intentional for causal proposal pairing,
but the two baselines answer different questions. The candidate-0 comparison
isolates EIG selection within a three-proposal interface; the thinking-naive
comparison estimates end-to-end policy value.

Finally, the method is one-step. Full two-step lookahead failed in the small
development study and is not rehabilitated here. MediQ transfer has not yet been
run, so no medical or cross-domain generality claim is made. A positive Paprika
result would support troubleshooting performance at a fixed turn budget, not a
universal claim that Bayesian arbitration improves every interactive agent.

\section{Conclusion}

Faithful BED transfer can fail when the proposal interface is detached from the
task objective. We test a narrower use of beliefs: let a thinking model propose
good actions, then use one-step EIG only to select among them when the estimated
advantage is large enough. Paprika provides an external sequential benchmark,
and a prompt-matched candidate-0 control isolates selection from elicitation.
The method and analysis are frozen before held-out outcomes are opened. Whether
beliefs improve troubleshooting selection is therefore an empirical question,
not a conclusion imported from the development pilot.

\bibliographystyle{plainnat}
\bibliography{references}

\end{document}
